\clearpage
\section{配比模型跨规模绝对 Loss 的适用边界诊断}
\label{sec:appendix-q1-transfer}
\setcounter{table}{0}
\setcounter{equation}{0}
\setcounter{figure}{0}

本附录检查未经规模校准的 1M 配比模型是否能够直接预测较大规模组的绝对 Loss，用于说明第一问的适用边界。该模型只以训练配比为输入；这一检查不评价第二问含规模项的标度模型，也不构成对最终四变量模型的验证。

表\ref{tab:q1-r2-transfer}列出冻结 1M 模型在三组比较数据上的指标。1M 同规模组的 13 个验证域均有正的 $R^2$；直接用于 60M、1B 时，两组的 13 个域均为负，表明未经校准的绝对数值不能跨规模直接使用。

\begin{table}[htbp]
 \centering\small
 \caption{冻结 1M 配比模型的绝对 Loss 迁移诊断（各指标为 13 域中位数）}
 \label{tab:q1-r2-transfer}
 \begin{tabular}{lrrrr}
 \toprule
 比较规模 & 配方数 & $R^2$ & RMSE & MAPE\\
 \midrule
 1M（同规模） & 256 & 0.6710 & 0.4227 & 6.95\%\\
 60M（直接迁移） & 256 & $-8.3331$ & 1.4270 & 37.40\%\\
 1B（直接迁移） & 64 & $-506.7317$ & 2.9429 & 147.39\%\\
 \bottomrule
 \end{tabular}
\end{table}

模型未校准规模变化引起的 Loss 水平和幅度变化。例如，Pile-CC 在 60M、1B 中的平均预测偏差分别为 1.0621、2.5206，而观测标准差仅为 0.3138、0.1008。系统偏差相对于组内波动较大，因而产生显著负的 $R^2$。较高的排序相关与此并不矛盾：排序要求相对次序接近，绝对预测还要求数值吻合。因此，第一问的跨规模比较仅检验配方排序；后续绝对 Loss 采用第二问的规模主干及明示连接假设，本表不能验证跨附件相对配比效应的幅度。
